\section{量子多体物理理论的发展}

\subsection{早期：量子统计与理想气}

量子力学建立后，玻色与费米统计立刻被用于理想气体与黑体辐射等问题。理想费米气给出金属电子气的基本图像：费米球、态密度、泡利顺磁等。理想模型虽忽略相互作用，却已能解释相当一部分低温金属行为，说明“费米面”是一个稳健概念。

\subsection{二次量子化与量子场论语言}

为处理可变粒子数与交换对称性，产生/湮灭算符与 Fock 空间成为标准语言。这一框架与量子场论同源：粒子作为场的量子激发出现，相互作用可用微扰级数与图展开处理。凝聚态中的准粒子、声子、磁振子等概念，都建立在这一语言之上。

\subsection{格林函数、图技术与费米液体}

战后，量子多体格林函数方法迅速发展。Matsubara（虚时）技巧使有限温度计算系统化；Feynman 图与 Dyson 方程给出自能与屏蔽等直观图像。Landau 费米液体理论表明：在三维金属中，相互作用往往把电子重整为带有有效质量和朗道参数的准粒子，低能激发仍可用“重整化费米气体”描述。

与此同时，人们也逐渐认识到费米液体并非普适：一维体系、莫特绝缘体、高温超导等要求超越简单准粒子图像。

\subsection{平均场、密度泛函与第一性原理}

Hartree--Fock 是最早的系统平均场近似之一，它包含交换作用但忽略动力学关联。Hohenberg--Kohn 与 Kohn--Sham 建立的密度泛函理论（DFT），把基态问题形式上映射为有效单体问题，使固体电子结构计算成为可能。局域密度近似（LDA）与广义梯度近似（GGA）等交换关联近似，在材料计算中取得了巨大成功。

\begin{note}[思政融入：理论联系实际与家国情怀]
密度泛函等方法已广泛服务于材料设计、催化、能源与信息器件等实际问题。中国学者在计算凝聚态、强关联与拓扑物性等领域作出了重要贡献。学习与应用这些理论时，既要追求学术深度，也要关注其服务国家重大需求的可能，发扬爱国主义精神与责任感。
\end{note}

\subsection{强关联、拓扑与当代前沿}

当相互作用与动能相当或更大时，平均场与弱耦合图技术可能失效。Hubbard 模型、Anderson 杂质模型、动力学平均场（DMFT）、张量网络与量子信息视角，成为当代多体物理的重要分支。此外，拓扑物态强调能带与纠缠的全局结构，与传统“准粒子 + 响应”图像互补。

本课程按大纲选取一条经典主线：\textbf{二次量子化 $\to$ 格林函数 $\to$ Hartree--Fock/DFT $\to$ 线性响应}。掌握这条主线后，再进入强关联或拓扑专题会更为顺畅。

\subsection{本课程学习建议}

\begin{enumerate}
  \item 熟练掌握二次量子化代数与常见哈密顿量的写法；
  \item 理解 retarded/advanced/Matsubara 格林函数及其谱表示；
  \item 能用 Wick 定理与简单 Feynman 图计算最低阶自能与响应；
  \item 弄清 Hartree、Fock、交换与关联的物理含义，以及 KS-DFT 的逻辑；
  \item 会从线性响应出发写出密度响应与 Lindhard 函数，并联系实验可观测量。
\end{enumerate}

考核方式为平时作业与综合项目。建议边学边做小计算（解析或数值），把抽象定义落到具体模型上。
